\chapter{Introducere}

\section{Scurtă prezentare a aplicației}
\vspace{0.8cm}
\begin{center}
\begin{figure}[htbp]
\includegraphics[width=0.72\textwidth]
{imagini/promonnou.eps}
\end{figure}
\end{center}
\paragraph{ }Trăim într-o societate în care una dintre activitățile zilnice întreprinse de o persoană este vânatoarea de promoții. ``Comorile'' zilelor noastre, promoțiile, ne ajută deseori să reducem suma cheltuită pe un produs care ne era necesar și astfel să ne permitem să achiziționăm produse chiar cu un buget scăzut.\newline 
După cum se știe, există și o ,,sărbătoare'' în acest scop, și anume, ,,Black Friday''\footnote{traducere ,,Vinerea Neagră''}, fiind extinsă deseori la o săptamână sau chiar mai mult, în care multe firme dețin o gamă largă de promoții, realizând în doar câteva zile, încasări record.\newpage Aceasta este însă de multe ori doar o strategie de marketing, iar produsele care par să aibă un preț foarte avantajos să fie de fapt unul normal sau chiar mai ridicat ca în alte perioade ale anului. De aceea este indicat să căutăm promoții atunci cand chiar avem nevoie de un produs si să urmarim variația prețului acestuia pe o perioadă de timp mai lungă. Promoțiile se desfășoară pe tot parcursul anului, chiar dacă nu sunt atat de mediatizate.\newline
 Există promoții la aproape orice fel de produs iar căutarea acestora cu ajutorul motorului de căutare Google ne poate furniza atât informații utile dar și informații redundante, produse ale căror denumire conține un cuvânt cheie atribuit atât  produsului căutat dar și altor produse care pot să nu aibă nicio legătură cu căutarea dvs. De aceea m-am gândit la implementarea unei aplicații care să grupeze laolaltă produsele promoționale de pe mai multe site-uri divizându-le în categorii sugestiv.\newline
 Consider că această aplicație este de mare ajutor în zilele noastre, optimizând atât procesul de căutare, reducând durata acestuia, având o interfață prietenoasă și ușor de utilizat de către oricine deține un smartphone sau o tabletă cu sistemul de operare Android.\bigskip

\section{De ce Android?}
\paragraph{ }
Am ales să dezvolt o aplicație pe un dispozitiv mobil, întrucât smartphone-urile și tabletele devin din ce  în ce mai populare și mai des utilizate, chiar dacă au apărut pe piața  mondială relativ recent, primul model în 1992, numit IBM Simon  \footnote{\burl{http://www.qrcodescanning.com/smartphonehist.html}}.\newline 
Conform datelor iSense.Solutions, platforma Android are un market share de peste 81\% fiind lider de piață atât în România cât și pe plan mondial.\newline
Exceptând Apple, care a ales iOS, și Microsoft, care utilizează sistemul de operare Windows Phone, aproape toate companiile producătoare de telefoane mobile folosesc sistemul de operare al celor de la Google, Android.\newline

Astfel, consumatorii au de ales dintr-o  arie foarte largă de telefoane având ca sistem de operare Android, și în acest fel pot să-și aleagă dintre toate modelele si dintre toți producătorii ceea ce li se potrivește cel mai bine. Platforma Android se regăsește de asemenea si pe tablete, ceasuri sau televizoare, numărul total de vânzari de smartphone-uri cu Android depășind 1 miliard în 2014. \footnote{\burl{http://www.businessinsider.com/android-1-billion-shipments-2014-strategy-analytics-2015-2}}\newpage
\begin{center}
\begin{figure}[htbp]
\caption{Android vs alte sisteme de operare}
\includegraphics[width=0.8\textwidth]
{imagini/numberAndroid.eps}
\end{figure} 
\end{center}

Unul dintre cele mai cunoscute limbaje de programare este Java\footnote{mai multe detalii în Capitolul 4 - Limbajul de programare Java}. Dezvoltat în anii 1990 de către Sun Microsystems, Java, la fel ca și platforma Android,  se bucură de o popularitate crescută,  aflându-se printre cele mai căutate limbaje de programare. Limbajul de programare Java este folosit atât pentru crearea conținutului paginilor web, aplicațiilor desktop, a jocurilor, precum și pentru aplicațiile pentru sistemul de operare Android.
Aplicațiile Java pot fi rulate pe mai multe platforme, de exemplu, un program scris în acest limbaj de programare poate fi folosit și pe sistemul de operare Windows și pe Mac OS X.\newline 

Am ales astfel să dezvolt o aplicație având o temă actuală, utilizând ca mediu de programare platforma Android deoarece aceasta este cea mai utilizată la nivel mondial dar și din dorința de a învăța lucruri noi și de a-mi dezvolta cunoștințele în această direcție având în vedere că tehnologia este în continuă dezvoltare iar eu ca programator trebuie să țin pasul cu ea.\newline

